\section{Introduction}


\revise{Software systems are increasingly used to manage real-world
assets of significant value. They hold currencies and the digital
wallets that store them, underpin the operations of banks and other
financial institutions, and keep much of the sensitive personal
information that was once recorded on paper. 
This trend is large and
accelerating: the IMF's 2025 Financial 
Access Survey reports that
digital transactions, including mobile 
money and mobile and internet
banking, have surged in emerging and 
developing economies, rising
from 55 transactions per adult in 2017 to 251 per adult by
2024~\cite{imf2025fas}. As more value 
moves into software, a single
defect or compromise can translate directly into 
real-world harm.}

\revise{Not only has software itself changed, but so 
has the way
software is built. Despite well-known
limitations of their black-box nature,
it is by now an
undeniable fact that developers 
routinely rely on AI-assisted coding
tools to help implement software. These 
tools dramatically lower the
barrier to development: as long as a 
need can be described in natural
language, the tool can try its best to 
help produce a 
corresponding implementation.
This opens software development to a far wider audience, including
people with little formal programming 
background. Tools such as Claude
Code~\cite{anthropic_claude_code}, Codex~\cite{openai_codex_2026}, and
OpenCode~\cite{opencode_2026} are now widely 
used in practice. }

\revise{Not only has the way software is built changed, but so has the
question of who, or what, uses it.
Traditionally, software was
designed with specific users in mind: when 
building an application, a
developer could reasonably assume it would 
be operated by humans, for
instance tapping on an iPhone screen. 
This assumption no longer holds.
As soon as software exposes 
programmable interfaces, those interfaces
can be invoked by autonomous 
agents in arbitrary ways that the
developers may never have 
anticipated at design time.}

\revise{The above two facts, despite the convenience and benefits they
bring, open new risks compared with the traditional workflow of
developing software. Figure~\ref{fig:workflow} further clarifies the
new risks introduced on top of the existing software development
lifecycle. First, in the software development stage, not only do humans
make mistakes in the limited number of lines they write themselves, but
errors can also arise in the thousands of lines of code generated by
AI-assisted coding tools. Such a large volume of code makes it
difficult even for expert developers to review, let alone the new
non-technical programmers who have not received sufficient software
training. Second, in the software deployment stage, when monitoring
users' activity, traditional metrics such as click rates and
day--night usage patterns no longer apply. Human users may only use
social media when they are awake, which follows a typical biological
cycle, but autonomous agents can access the software in unprecedented
ways. This new usage, often deviating from developers' original intent,
can pose serious risks to software that stores critical real-world
value, such as digital currency, credentials, and classified
blueprints.}

\begin{figure}[t]
  \centering
  \includegraphics[width=\textwidth]{figures/workflow.pdf}
  \caption{\revise{The overall workflow for developing
  high-stakes software.
  In \emph{Stage~1 (development)}, software is written by
  human developers together with AI-assisted coding tools, which
  introduce new risks into the code itself. After review, testing,
  auditing, and verification, the software moves to \emph{Stage~2
  (deployment)}, where it is invoked not only by human users but
  also by autonomous agents, introducing further risks.}}
  \label{fig:workflow}
\end{figure}

\revise{Because of these shifts, we are already witnessing rapidly
growing security threats to software. Cyberattacks are not a new
problem. The IMF reports that the financial sector has lost \$12
billion over the past two decades to more than 20{,}000 such
incidents~\cite{financemagnates2024imf}. However, the two changes
above sharply amplify the threat. On the development side,
AI-generated code has already caused direct financial loss. In early
2026, a faulty price-oracle update in a pull request co-authored by an
AI coding model mispriced collateral on the Moonwell lending protocol,
leading to roughly \$1.78 million in losses~\cite{moonwell2026exploit}.
On the usage side, autonomous agents have themselves been turned into
attack tools. A state-sponsored group weaponized an AI coding agent to
run a largely autonomous cyber-espionage campaign against around
thirty organizations before the provider detected the activity and
disabled the accounts~\cite{anthropic2025espionage}. These attacks
compound an already worsening landscape. In 2025, a record year for
cryptocurrency theft, attackers stole more than \$3 billion in digital
assets~\cite{chainalysis2025crypto}.}

\revise{It is therefore increasingly important to protect what we call
\emph{high-stakes software}. Informally, a high-stakes software
system is durable software whose purpose is to protect
and operate on critical assets such as digital wallets, secrets,
credentials, and confidential documents. Such systems naturally
attract attackers, and we believe they are the most valuable to
defend. In contrast, disposable, single-use programs that are written
once and discarded fall outside our scope. This thesis focuses on
securing high-stakes software against the two untrusted factors
introduced above: AI-assisted coding tools at development time, and
autonomous agents at deployment time. We concretize this goal in two
domains: the security of Web3 software, and the mitigation of harmful
behavior in LLM-based systems.}




\subsection{High-Stakes Software}

A software system is \emph{high-stakes} when its execution can materially affect
assets or decisions that matter outside the software itself, such as money,
secrets, credentials, personal data, operational infrastructure, or access
rights. The term is not tied to any technical domain. A banking backend, a
medical-control system, a confidential-document platform, a smart contract, and
even a short generated script that is handed a private key can all be high-stakes
once they are entrusted with critical assets. What makes software high-stakes is
the consequence of its failure, not the size or complexity of its code. A small
program can be high-stakes if it holds a secret, while a large application may be
low-stakes if failure only causes inconvenience; a complex text parser, for
instance, is not high-stakes when its mistakes are merely an annoyance.

\begin{definition}[High-stakes software]
\label{def:high-stakes}
A high-stakes software is a software artifact or deployed system that mediates
access to critical assets and whose violation of intended security properties
can cause material harm to users, organizations, or the public.
\end{definition}

This definition covers both long-lived deployed systems and the software
artifacts that become part of them. We focus on durable systems and repeated
workflows, where repeated use gives attackers more opportunity and makes
defensive investment more worthwhile. The broader principle is that whenever
software is trusted with a critical asset, that trust must be justified. This
focus is increasingly important: as code becomes cheap to produce and easy to
generate with language models, security effort is best concentrated on the
software that actually guards critical assets, rather than on disposable code
such as one-off parsers whose failure carries no real cost.



\subsection{Untrusted Development}

High-stakes software is not built in isolation. It is assembled from many
influences: human programmers, generated code, libraries and packages,
documentation and online examples, configuration, build systems, and review
tools. Development becomes \emph{untrusted} when some of these influences cannot
be assumed correct, benign, or aligned with the developer's intent. This does
not require any party to be malicious; it only requires that the safety of an
influence cannot be fully established before its output is accepted. A code
generator may fabricate an API, a documentation page may be stale or poisoned, a
dependency may hide unwanted behavior, a review tool may miss a subtle bug, or a
developer may accept a patch that looks correct locally yet breaks a global
security assumption. In each case, the software can carry harmful behavior before
it is ever deployed.

\begin{definition}[Untrusted development]
\label{def:untrusted-dev}
Development is untrusted when the construction, selection, modification, or
review of software depends on actors, tools, data sources, or artifacts whose
correctness, provenance, or alignment cannot be fully verified at the time they
influence the software.
\end{definition}

The defining feature of untrusted development is that the risk is embedded in the
artifact: once a harmful dependency, patch, or code fragment is accepted, the
deployed system will faithfully execute it. Development-time defenses therefore
aim to detect, reject, constrain, or harden such artifacts before they become
trusted software.



\subsection{Untrusted Use}

Even carefully written, reviewed, and tested software remains only as secure as
the way it is used. Developers design software with an expected usage model in
mind: who will call the system, in what order, at what scale, through which
components, and what counts as normal behavior. These assumptions are usually
implicit, and they may hold for ordinary users yet fail once the system meets
adversarial users, automated clients, composed services, or autonomous agents.
Use becomes \emph{untrusted} when a deployed system can be invoked by parties
whose goals, timing, scale, identity, or interaction patterns cannot be assumed
to match the developer's intent. Such misuse need not exploit a traditional bug
in any single function. A harmful execution may consist entirely of permitted
operations arranged in an unexpected order, at an unusual scale, through an
unusual caller, or in combination with other systems. The Harvest example in
Chapter~1 illustrates this: the exploit transaction was abnormal in its
invocation pattern, repetition, gas consumption, and asset movement, even though
each of its individual calls was allowed by the code.

\begin{definition}[Untrusted use]
\label{def:untrusted-use}
Use is untrusted when a deployed software system is invoked through inputs,
users, agents, programs, or interaction patterns whose behavior cannot be
assumed to be benign, human-like, anticipated, or aligned with the system
designer's intent.
\end{definition}

The defining feature of untrusted use is that the risk is induced at execution
time. The artifact may be unchanged and every individual operation legal, yet
the resulting execution trace can still violate the assumptions under which the
system was designed. Use-time defenses therefore aim to monitor, characterize,
restrict, or reject executions that fall outside acceptable behavior.

Together, untrusted development and untrusted use are two forms of what we call
\emph{untrusted influence} over high-stakes software: any actor, artifact, data
source, tool, or interaction pattern that shapes the behavior of a high-stakes
system without a sufficient basis for trust. Development-time influence shapes
what the software is; use-time influence shapes what the software does. The
security goal of this dissertation is to limit how much untrusted influence can
affect critical assets.

More formally, let $S$ be a software system that mediates a set of critical
assets $A$, and let $\Phi(A)$ denote the intended security properties of those
assets, such as confidentiality, integrity, authorization, availability, or
conservation of value. The behavior of $S$ is shaped in two phases: during
development, a collection of influences $\mathcal{I}_{\mathrm{dev}}$ produces the
software artifact, and during deployment, a collection of influences
$\mathcal{I}_{\mathrm{use}}$ produces executions of that artifact. The system is
under untrusted influence when some element of $\mathcal{I}_{\mathrm{dev}} \cup
\mathcal{I}_{\mathrm{use}}$ can affect whether $\Phi(A)$ holds yet cannot be
assumed correct, benign, or aligned with the system's intended behavior. This
dissertation studies defenses at both points in the lifecycle: before
deployment, by auditing and hardening artifacts produced under untrusted
development, and after deployment, by detecting or blocking executions produced
under untrusted use.





